\documentclass{article}
\usepackage[T1]{fontenc}
\usepackage{lmodern}
\usepackage{graphicx}
\usepackage{amsmath,amssymb}
\usepackage[a4paper,margin=2cm]{geometry}
\usepackage[absolute,overlay]{textpos}
\setlength{\TPHorizModule}{1mm}
\setlength{\TPVertModule}{1mm}
\usepackage[indonesian]{babel}
\usepackage{hyperref}
\setlength{\emergencystretch}{2em}
% unit: Halaman 13

\begin{document}

% === Halaman 13 (llm) ===

\begin{itemize}
  \item Contoh pada Caesar Cipher, penyerang memilih plainteks sederhana, misalnya:
  
  AAAAA
  
  lalu mengirim plainteks ke sistem enkripsi
  
  Sistem mengembalikan cipherteks, misalnya:
  
  FFFFF
  
  Dari sini terlihat huruf A $\rightarrow$ F, berarti pergeseran = 5.
\end{itemize}

Dengan kunci itu, penyerang bisa mendekripsi semua cipherteks lain.

\end{document}
